\uuid{485v}
\exo7id{1671}
\titre{exo7 1671}
\auteur{roussel}
\organisation{exo7}
\datecreate{2001-09-01}
\isIndication{false}
\isCorrection{false}
\chapitre{Réduction d'endomorphisme, polynôme annulateur}
\sousChapitre{Diagonalisation}
\cpge{mpsi-mp2i.calcul-matriciel-systemes.e, pcsi.determinants.c, ptsi.determinants.a, mpsi-mp2i.matrices-applications-lineaires.b, pcsi.matrices-applications-lineaires.b, ptsi.matrices-applications-lineaires.b, bcpst1.matrices, bcpst1.suites-reelles, mp-mpi.reduction-endomorphismes.d, pc.reduction.c, psi.reduction-endomorphismes.c, pt.reduction-endomorphismes.c, mp-mpi.reduction-endomorphismes.i, pc.reduction.d, psi.reduction-endomorphismes.d, bcpst2.valeurs-propres-vecteurs-propres.b}
\module{Algèbre}
\niveau{L2}
\difficulte{}

\contenu{
\texte{
Soit $B$ une matrice diagonalisable ~de $\mathcal{M}_n \left (\mathbb{R} \right )$.
On d\'efinit son rayon spectral par
$$\rho(B)= \max \left \{ |\lambda| \mbox{~~avec~~} \lambda \mbox{~~est une
valeur propre ~de~~} B \right \}.$$
}
\begin{enumerate}
    \item \question{Montrer que $\lim _{k \longrightarrow + \infty} B^k= 0$.}
    \item \question{En d\'eduire que $I-B$ est inversible et que $(I-B)^{-1}=
\displaystyle{\sum _{k=0}^{+ \infty}} B^{k}$.}
\end{enumerate}
}
